\documentclass[a4paper,12pt]{article}

\usepackage[T2A]{fontenc}
\usepackage[utf8]{inputenc}
\usepackage[russian]{babel}
\usepackage{amsmath,amssymb}
\usepackage[scaled=0.96]{PTSans}
\renewcommand{\familydefault}{\sfdefault}
\usepackage{icomma}
\usepackage[a4paper,margin=22mm,headheight=25pt]{geometry}
\usepackage{microtype}
\usepackage{tikz}
\usetikzlibrary{calc,arrows.meta,positioning,shapes.geometric}
\usepackage{pgfplots}
\pgfplotsset{compat=1.18}
\usepackage{tabularx,booktabs}
\usepackage{enumitem}
\usepackage{parskip}
\usepackage{fancyhdr}
\usepackage{setspace}
\usepackage{xcolor}

\definecolor{accent}{HTML}{008080}
\definecolor{darkaccent}{HTML}{004D4D}
\definecolor{textgray}{HTML}{2F3333}
\definecolor{softgray}{HTML}{E7ECEC}
\definecolor{midgray}{HTML}{6B7474}

\color{textgray}
\onehalfspacing
\setlength{\parindent}{0pt}
\setlength{\emergencystretch}{2em}

\setlist[itemize]{
    leftmargin=1.6em,
    itemsep=0.30em,
    topsep=0.30em
}

\setlist[enumerate]{
    leftmargin=1.9em,
    itemsep=0.40em,
    topsep=0.35em
}

\pagestyle{fancy}
\fancyhf{}
\fancyhead[L]{\small Графики функций: гипербола}
\fancyhead[R]{\small Марина Селиверстова}
\fancyfoot[C]{\small\thepage}
\renewcommand{\headrulewidth}{0.4pt}
\renewcommand{\headrule}{%
    \hbox to\headwidth{%
        \color{softgray}%
        \leaders\hrule height \headrulewidth\hfill
    }%
}

\newcommand{\sectionrule}{%
    \vspace{-0.4em}%
    {\color{softgray}\hrule}%
    \vspace{0.7em}%
}

\newcommand{\term}[1]{\textcolor{darkaccent}{\textbf{#1}}}
\newcommand{\R}{\mathbb{R}}

\tikzset{
    startstop/.style={
        ellipse,
        draw=darkaccent,
        very thick,
        align=center,
        minimum width=38mm,
        minimum height=10mm,
        fill=softgray,
        font=\small
    },
    process/.style={
        rounded corners=2mm,
        rectangle,
        draw=accent,
        thick,
        align=center,
        text width=71mm,
        minimum height=11mm,
        inner sep=3mm,
        font=\small
    },
    decision/.style={
        diamond,
        aspect=2.35,
        draw=darkaccent,
        thick,
        align=center,
        text width=39mm,
        inner sep=1.5mm,
        font=\small
    },
    branchprocess/.style={
        rounded corners=2mm,
        rectangle,
        draw=accent,
        thick,
        align=center,
        text width=50mm,
        minimum height=13mm,
        inner sep=2.5mm,
        font=\small
    },
    flowarrow/.style={
        -{Stealth[length=3mm,width=2mm]},
        thick,
        draw=darkaccent
    },
    flowlabel/.style={
        font=\small,
        fill=white,
        inner sep=1pt
    }
}

\begin{document}

% =========================================================
% ТИТУЛЬНЫЙ ЛИСТ
% =========================================================

\begin{titlepage}
\thispagestyle{empty}
\centering

\vspace*{23mm}

{\Large\textcolor{darkaccent}{\textbf{Чек-лист}}\par}

\vspace{6mm}

{\huge\bfseries
Гипербола и обратная пропорциональность
\par}

\vspace{3mm}

{\Large\bfseries
Асимптоты, сдвиги и выколотые точки
\par}

\vspace{18mm}

{\large Марина Селиверстова\par}
\vspace{2mm}
{\large ОГЭ\par}

\vfill

{\large Лёвин Артём Александрович\par}
\vspace{2mm}
{\normalsize эксклюзивно для Марины Селиверстовой\par}
\vspace{5mm}
{\normalsize \today\par}

\vspace{21mm}

\begin{tikzpicture}[remember picture,overlay]
\draw[
    accent,
    line width=1.3pt,
    opacity=0.55
]
    ($(current page.south west)+(0,16mm)$)
    .. controls
    ($(current page.south west)+(55mm,28mm)$)
    and
    ($(current page.south east)+(-60mm,5mm)$)
    ..
    ($(current page.south east)+(0,18mm)$);
\end{tikzpicture}

\end{titlepage}

% =========================================================
% 1. БАЗОВАЯ ГИПЕРБОЛА
% =========================================================

\section*{1. Базовая гипербола $y=\dfrac{k}{x}$}
\sectionrule

\term{Обратная пропорциональность} задаётся формулой
\[
y=\frac{k}{x},
\qquad k\ne0.
\]
Её график --- \term{гипербола}, состоящая из двух ветвей.

\begin{itemize}
    \item \textbf{Область определения:}
    \[
    x\in\R\setminus\{0\}.
    \]

    \item \textbf{Множество значений:}
    \[
    y\in\R\setminus\{0\}.
    \]

    \item \textbf{Асимптоты:}
    \[
    x=0,
    \qquad
    y=0.
    \]
    График не имеет с ними общих точек.

    \item Если $k>0$, ветви расположены в I и III четвертях.

    \item Если $k<0$, ветви расположены во II и IV четвертях.

    \item Для любой точки гиперболы выполняется
    \[
    xy=k.
    \]
    При фиксированном $|x|$ увеличение $|k|$ увеличивает $|y|$;
    при фиксированном $|y|$ увеличивается и $|x|$.
\end{itemize}

\subsection*{Как построить по точкам}

Выберите несколько ненулевых значений $x$ по обе стороны от нуля,
вычислите $y=\dfrac{k}{x}$, нанесите точки и проведите две плавные ветви,
которые приближаются к асимптотам.

Для $y=\dfrac1x$ удобно взять, например:

\begin{table}[ht]
\centering
\renewcommand{\arraystretch}{1.25}
\begin{tabular}{@{}ccccccc@{}}
\toprule
$x$ & $-2$ & $-1$ & $-\dfrac12$ & $\dfrac12$ & $1$ & $2$ \\
\midrule
$y=\dfrac1x$ & $-\dfrac12$ & $-1$ & $-2$ & $2$ & $1$ & $\dfrac12$ \\
\bottomrule
\end{tabular}
\end{table}

\begin{center}
\begin{tikzpicture}
\begin{axis}[
    width=0.86\linewidth,
    height=8.8cm,
    axis lines=middle,
    xmin=-5,
    xmax=5,
    ymin=-5,
    ymax=5,
    xtick={-4,-3,-2,-1,1,2,3,4},
    ytick={-4,-3,-2,-1,1,2,3,4},
    samples=220,
    unit vector ratio=1 1,
    xlabel={$x$},
    ylabel={$y$},
    legend style={
        at={(0.98,0.98)},
        anchor=north east,
        draw=none,
        fill=white,
        font=\small
    }
]

\addplot[accent,thick,domain=-5:-0.21,restrict y to domain=-5:5]{1/x};
\addlegendentry{$y=\dfrac1x$}
\addplot[accent,thick,domain=0.21:5,restrict y to domain=-5:5,forget plot]{1/x};

\addplot[darkaccent,dashed,thick,domain=-5:-0.41,restrict y to domain=-5:5]{2/x};
\addlegendentry{$y=\dfrac2x$}
\addplot[darkaccent,dashed,thick,domain=0.41:5,restrict y to domain=-5:5,forget plot]{2/x};

\addplot[midgray,dotted,very thick,domain=-5:-0.11,restrict y to domain=-5:5]{0.5/x};
\addlegendentry{$y=\dfrac{0{,}5}{x}$}
\addplot[midgray,dotted,very thick,domain=0.11:5,restrict y to domain=-5:5,forget plot]{0.5/x};

\end{axis}
\end{tikzpicture}
\end{center}

\textbf{Самопроверка.}
Перед построением определите знак $k$ и четверти, затем убедитесь,
что ни одна ветвь не пересекает оси координат.

% =========================================================
% 2. СДВИНУТАЯ ГИПЕРБОЛА
% =========================================================

\section*{2. Линейно-дробная функция и асимптоты}
\sectionrule

Рассмотрим функцию
\[
y=\frac{ax+b}{cx+d},
\qquad c\ne0.
\]
Если
\[
ad-bc\ne0,
\]
то её график является сдвинутой гиперболой.

\textbf{Сначала обязательно учитываем ограничение}
\[
cx+d\ne0,
\qquad
x\ne-\frac{d}{c}.
\]

Асимптоты имеют вид
\[
\boxed{x=-\frac{d}{c}},
\qquad
\boxed{y=\frac{a}{c}}.
\]

Это удобно увидеть из преобразования
\[
\frac{ax+b}{cx+d}
=
\frac{a}{c}
+
\frac{bc-ad}{c(cx+d)}.
\]
Точка пересечения асимптот
\[
\left(-\frac{d}{c},\frac{a}{c}\right)
\]
является центром гиперболы.

\subsection*{Пример}

Пусть
\[
y=\frac{2x-3}{x+5}.
\]
Область определения:
\[
x\ne-5.
\]
Преобразуем:
\[
\frac{2x-3}{x+5}
=
2-\frac{13}{x+5}.
\]
Поэтому сразу видны асимптоты
\[
\boxed{x=-5},
\qquad
\boxed{y=2}.
\]

Для пересечения с осью $Ox$ полагаем $y=0$:
\[
2x-3=0
\quad\Longrightarrow\quad
x=\frac32.
\]
Точка пересечения:
\[
\left(\frac32,0\right).
\]

Для пересечения с осью $Oy$ подставляем $x=0$:
\[
y=-\frac35.
\]
Точка пересечения:
\[
\left(0,-\frac35\right).
\]

\begin{center}
\begin{tikzpicture}
\begin{axis}[
    width=0.82\linewidth,
    height=8.2cm,
    axis lines=middle,
    xmin=-10,
    xmax=4,
    ymin=-5,
    ymax=7,
    xtick={-9,-7,-5,-3,-1,1,3},
    ytick={-4,-2,2,4,6},
    samples=220,
    xlabel={$x$},
    ylabel={$y$},
    unit vector ratio=1 1,
    legend style={draw=none,fill=white,font=\small,at={(0.98,0.97)},anchor=north east}
]
\addplot[accent,thick,domain=-10:-5.12,restrict y to domain=-5:7]{(2*x-3)/(x+5)};
\addlegendentry{$y=\dfrac{2x-3}{x+5}$}
\addplot[accent,thick,domain=-4.88:4,restrict y to domain=-5:7,forget plot]{(2*x-3)/(x+5)};
\addplot[darkaccent,dashed,thick] coordinates {(-5,-5) (-5,7)};
\addplot[darkaccent,dashed,thick] coordinates {(-10,2) (4,2)};
\addplot[only marks,mark=*,mark size=2pt] coordinates {(1.5,0) (0,-0.6)};
\end{axis}
\end{tikzpicture}
\end{center}

\textbf{Практический приём.}
Для эскиза обычно достаточно найти асимптоты и несколько надёжных точек,
в частности пересечения с осями, если они существуют.

% =========================================================
% 3. СОКРАЩЕНИЕ И ВЫКОЛОТАЯ ТОЧКА
% =========================================================

\section*{3. Сокращение дроби: исходные ограничения сохраняются}
\sectionrule

Рассмотрим
\[
f(x)=\frac{2x+1}{2x^2+x}
=\frac{2x+1}{x(2x+1)}.
\]
Сначала учитываем \textbf{исходный} знаменатель:
\[
x(2x+1)\ne0
\quad\Longrightarrow\quad
x\ne0,
\qquad
x\ne-\frac12.
\]

После сокращения
\[
f(x)=\frac1x,
\qquad
x\ne0,-\frac12.
\]
Запрещённое значение $x=-\dfrac12$ исчезло из упрощённой формулы,
но осталось в области определения. По формуле $y=\dfrac1x$ ему соответствовало бы
$y=-2$, поэтому точка
\[
\left(-\frac12,-2\right)
\]
\term{выколота}.

\textbf{Правило.}
Сокращение упрощает формулу функции, но не возвращает значения $x$,
запрещённые исходным знаменателем.

\subsection*{Параметрическая прямая $y=kx$}

Прямая $y=kx$ проходит через начало координат. Для пересечения с
$y=\dfrac1x$ имеем
\[
kx=\frac1x
\quad\Longrightarrow\quad
kx^2=1,
\qquad x\ne0.
\]
При $k>0$ у обычной гиперболы две точки пересечения.
Чтобы у рассматриваемой функции осталась ровно одна общая точка,
одна из них должна совпасть с выколотой точкой
$\left(-\dfrac12,-2\right)$.

Подставим её в $y=kx$:
\[
-2=k\left(-\frac12\right)
\quad\Longrightarrow\quad
\boxed{k=4}.
\]
Проверка:
\[
4x^2=1
\quad\Longrightarrow\quad
x=\pm\frac12.
\]
Значение $x=-\dfrac12$ запрещено, а $x=\dfrac12$ допустимо.
Следовательно, при $k=4$ остаётся ровно одна общая точка.

% =========================================================
% БЛОК-СХЕМА
% =========================================================

\clearpage
\newgeometry{left=18mm,right=18mm,top=15mm,bottom=15mm}
\thispagestyle{empty}

\begin{center}
{\Large\bfseries\textcolor{darkaccent}{Алгоритм построения рациональной функции гиперболического типа}}

\vspace{4mm}

\begin{tikzpicture}[node distance=4mm]

\node[startstop] (start) {Начало};

\node[process,below=of start] (odz)
{Выпишите ОДЗ по\\\textbf{исходному} знаменателю};

\node[process,below=of odz] (factor)
{Разложите числитель и знаменатель\\на множители};

\node[decision,below=6mm of factor] (common)
{Есть общий множитель?};

\node[branchprocess,below=9mm of common,xshift=-39mm] (cancel)
{Сократите его, сохранив исходные ограничения};

\node[branchprocess,below=9mm of common,xshift=39mm] (nocancel)
{Переходите к анализу\\без сокращения};

\node[process,below=39mm of common] (form)
{Приведите функцию к удобному виду\\и найдите асимптоты};

\node[process,below=of form] (holes)
{Проверьте запрещённые значения:\\после сокращения они могут дать\\выколотые точки};

\node[process,below=of holes] (points)
{Найдите удобные точки:\\пересечения с осями или значения из таблицы};

\node[process,below=of points] (sketch)
{Постройте ветви с учётом асимптот,\\ОДЗ и выколотых точек};

\node[startstop,below=of sketch] (finish) {Эскиз готов};

\draw[flowarrow] (start) -- (odz);
\draw[flowarrow] (odz) -- (factor);
\draw[flowarrow] (factor) -- (common);
\draw[flowarrow] (common) -- node[flowlabel,above left=1mm] {да} (cancel);
\draw[flowarrow] (common) -- node[flowlabel,above right=1mm] {нет} (nocancel);
\draw[flowarrow] (cancel) |- (form);
\draw[flowarrow] (nocancel) |- (form);
\draw[flowarrow] (form) -- (holes);
\draw[flowarrow] (holes) -- (points);
\draw[flowarrow] (points) -- (sketch);
\draw[flowarrow] (sketch) -- (finish);

\end{tikzpicture}
\end{center}

\clearpage
\restoregeometry

% =========================================================
% 4. ГОРИЗОНТАЛЬНЫЕ ПРЯМЫЕ
% =========================================================

\section*{4. Горизонтальная прямая $y=m$}
\sectionrule

Рассмотрим
\[
g(x)=3-\frac{x+2}{x^2+2x}.
\]
Исходный знаменатель даёт
\[
x(x+2)\ne0
\quad\Longrightarrow\quad
x\ne0,-2.
\]
При допустимых $x$ имеем $g(x)=3-\dfrac1x$.

Горизонтальная асимптота $y=3$, поэтому при $m=3$ общих точек нет.
Запрещённому значению $x=-2$ по упрощённой формуле соответствует
\[
y=3-\frac1{-2}=\frac72,
\]
поэтому точка $\left(-2,\dfrac72\right)$ выколота и при $m=\dfrac72$
пересечение тоже отсутствует.

Если $m\ne3$, то
\[
m=3-\frac1x
\quad\Longrightarrow\quad
x=-\frac1{m-3}.
\]
Это единственное решение; оно запрещено только при $x=-2$, то есть при
$m=\dfrac72$. Следовательно,
\[
\boxed{m=3\quad\text{или}\quad m=\frac72}.
\]

\section*{5. Мини-памятка и частые ошибки}
\sectionrule

\begin{itemize}
    \item Начинайте с ОДЗ \textbf{исходной} функции; после сокращения запрещённые значения сохраняются.
    \item Для $y=\dfrac{k}{x}$ оси координат --- асимптоты; ветви к ним приближаются, но не пересекают их.
    \item Знак $k$ определяет четверти: $k>0$ --- I и III; $k<0$ --- II и IV.
    \item Для $y=\dfrac{ax+b}{cx+d}$ при $c\ne0$ и $ad-bc\ne0$:
    $x=-\dfrac{d}{c}$ и $y=\dfrac{a}{c}$ --- асимптоты.
    \item Пересечение с $Ox$ ищут из $y=0$; с $Oy$ --- подстановкой $x=0$, если это допускает ОДЗ.
    \item Не путайте $y=kx$ и $y=m$: первая прямая проходит через начало координат, вторая горизонтальна.
\end{itemize}

% =========================================================
% ТРЕНИРОВОЧНЫЙ БЛОК
% =========================================================

\clearpage

\section*{Тренировочный блок --- Путь А}
\sectionrule

\textbf{1. Базовая гипербола}

Для функции
\[
y=-\frac3x
\]
выполните задания:

\begin{enumerate}[label=\alph*)]
    \item укажите область определения;
    \item назовите обе асимптоты;
    \item укажите четверти, в которых расположены ветви;
    \item заполните таблицу для $x=-3,-1,1,3$ и постройте эскиз графика.
\end{enumerate}

\vspace{7mm}

\textbf{2. Сокращение и одна общая точка}

Постройте график функции
\[
f(x)=\frac{3x-3}{x^2-x}
\]
с учётом всех ограничений.

Найдите значение параметра $a$, при котором прямая
\[
y=ax
\]
имеет с графиком функции ровно одну общую точку.

\vspace{7mm}

\textbf{3. Сдвиг, выколотая точка и горизонтальная прямая}

Постройте график функции
\[
g(x)=4-\frac{x-3}{x^2-3x}
\]
с учётом всех ограничений.

Найдите все значения $m$, при которых прямая
\[
y=m
\]
не имеет с графиком функции ни одной общей точки.

% =========================================================
% ОТВЕТЫ
% =========================================================

\clearpage

\section*{Ответы}
\sectionrule

\renewcommand{\arraystretch}{1.35}

\begin{tabularx}{\linewidth}{@{}>{\bfseries}p{0.08\linewidth}X@{}}
\toprule
№ & Ответ \\
\midrule
1 &
$D(y)=\R\setminus\{0\}$; асимптоты $x=0$ и $y=0$; ветви расположены во II и IV четвертях.
Для $x=-3,-1,1,3$ получаем соответственно $y=1,3,-3,-1$. \\
\addlinespace
2 &
$x\ne0,1$; после сокращения $f(x)=\dfrac3x$ при $x\ne0,1$; точка $(1,3)$ выколота.
Искомое значение:
\[
\boxed{a=3}.
\] \\
\addlinespace
3 &
$x\ne0,3$; после сокращения $g(x)=4-\dfrac1x$ при $x\ne0,3$; точка
$\left(3,\dfrac{11}{3}\right)$ выколота.
Искомые значения:
\[
\boxed{m=4\quad\text{или}\quad m=\frac{11}{3}}.
\] \\
\bottomrule
\end{tabularx}

\end{document}